\section{Speed, order and a measurement}\label{to:sec:measure}

\paragraph{Speed and order together.} Output deadlines alone impose no deadline on preparation or cleanup: a device can prepare its resources in advance, release its outputs on time and restore itself slowly afterwards. In our model the order cost is therefore pure order. A speed price appears only when the thermal interface is budgeted. Consider a collision model in which each thermal contact fully swaps one cell with a fresh Gibbs qubit. Let a device implement the target at error $\varepsilon$, and let it reset its cells along strokes of thermodynamic length $\ell_j$, each a cooling stroke that starts at excited population at most $1/2$, using $K_j$ contacts per stroke and no other operation on the cell during a stroke. Each contact is a step equilibration, and for a process carried out in $K$ such steps the dissipation is at least $\kB T\ell^2/(2K)$ to leading order in $1/K$, for thermodynamic length $\ell$~\cite{SalamonBerry83,Nulton85,Crooks07,ScandiPerarnauLlobet19,Abiuso20}. For these strokes the bound holds exactly at every $K_j$, as shown below, and summing over strokes gives
\[
W\ \ge\ W^{\mathrm{qs}}(\varepsilon)+\frac{\kB T}{2}\sum_j\frac{\ell_j^2}{K_j}\ \ge\ W^{\mathrm{qs}}(\varepsilon)+\kB T\,\frac{(\sum_j\ell_j)^2}{2K_{\mathrm{tot}}},
\]
where $W^{\mathrm{qs}}(\varepsilon)$ is the slow limit of the same device. For a device that attains its schedule's optimum it tends to $W^{\mathrm{qs}}_{\mathrm{batch}}+W_{\mathrm{order}}$ as $\varepsilon\to0$; any other device adds its own implementation overhead. For other strokes, including recoveries that extract work, the factor $1/2$ holds only to leading order in $1/K_j$; the factor $1/4$ holds exactly. Both exact statements have short proofs. A stroke whose contacts take the cell through the excited populations $a=q_0,q_1,\dots,q_K=b$ dissipates $\kB T\sum_jd(q_{j-1}\|q_j)$, with $d$ the binary relative entropy in nats, and $\ell=2|\theta(a)-\theta(b)|$ with $\theta(q)=\arcsin\sqrt q$. Jensen's inequality gives $d(p\|q)\ge-2\ln\cos(\theta(p)-\theta(q))\ge(\theta(p)-\theta(q))^2$ for every step, and Cauchy--Schwarz over the $K$ steps gives the factor $1/4$. For $q<p\le1/2$ the difference $d(p\|q)-2(\theta(p)-\theta(q))^2$ vanishes at $q=p$ and has derivative $[2(\theta(p)-\theta(q))-(p-q)/\sqrt{q(1-q)}]/\sqrt{q(1-q)}\le0$ in $q$, since $2(\theta(p)-\theta(q))=\int_q^p dx/\sqrt{x(1-x)}$ and $x(1-x)$ increases on $[0,1/2]$; so a downward step within $[0,1/2]$ costs at least $2(\theta(p)-\theta(q))^2$. In a cooling stroke from $a\le1/2$ to $b$, raising every population to at least $b$ and then replacing the path by its running minimum raises no step's cost and leaves only downward steps within $[b,a]$, and Cauchy--Schwarz gives the factor $1/2$. When ordinary and order-forced resets, of total lengths $L_B$ and $L_O$, share one budget of $M$ contacts, the least dissipation is $\kB T(L_B+L_O)^2/(2M)$ to leading order in $1/M$, with the contacts split in the ratio $L_B:L_O$. Two separate interfaces with the same total of $M$ contacts reach the same value when they split them in that ratio, so sharing an interface adds no cross term of its own; the difference $\kB T\,L_BL_O/M$ from $\kB T(L_B^2+L_O^2)/(2M)$ arises only against separate interfaces granted $M$ contacts each, $2M$ in all. The order-forced resets lengthen the path the contact budget must cover. Order decides which resets must happen; speed charges for completing them in time. At a fixed number of contacts, driving the error to zero makes the work diverge, as finite resources forbid a perfect reset~\cite{ReebWolf14,Taranto23}, although only logarithmically for a single contact and far more slowly for several. Any finite law must therefore state its error.

\paragraph{A measurement.} Measure the work as a function of the thermal-contact budget $K$ for both schedules at the same certified error. Fit each curve by an intercept plus terms in $1/K$ and $1/K^2$, or by the exact finite-$K$ curve of the collision model, as tests of Landauer's principle extrapolate the work to slow operation~\cite{Berut12,Jun14}. The difference between the two slow-limit intercepts is the excess quasistatic work of the two implementations at the certified error. It equals the order surcharge only when both implementations attain their schedule optima; otherwise it also contains implementation overhead, and the finite certificates below give an unavoidable lower bound without any such optimality claim. The slopes carry speed. Near a pure target the work approaches its $1/K$ law slowly: in the collision model above, for the amplitude-damping example below with a cell reset from its Gibbs state to excited population $10^{-6}$, a fit with $1/K$ alone over budgets of four to twelve contacts biases the intercept by about a tenth of the order cost, and adding the $1/K^2$ term, or fitting over ten to fifty contacts, brings the bias to about one percent. For amplitude damping at $p=1/2$ and $\beta E=\ln2$, the difference for many uses at vanishing error is $0.4313\,E$ per use, about $1.4\times10^{-24}$\,J at $5$\,GHz with a local reservoir at $346$\,mK. Over a finite number of uses it is smaller, and for a single use it vanishes.

\paragraph{Cautions.} Three cautions apply. First, at that temperature a $5$\,GHz bath implements generalized amplitude damping, so the target channel needs an engineered reservoir whose cost is counted. At impurity $10^{-6}$ the last thermal carrier needs a gap near $100$\,GHz, unless an encoded carrier or a costed colder reservoir is used. Second, amplitude damping is extreme, so its batch saving exists only at nonzero error, and both intercepts depend on that error. Third, the decisive comparison needs finite certificates for the implemented devices, not asymptotic rates.

\paragraph{Finite certificates.} Two finite certificates exist. One is a saving at $256$ uses and error $10^{-8}$ (Appendix~\ref{app:channels}); it needs reservoirs far too large to build. The other is the two-qubit process of the introduction, and it does not.

\begin{proposition}[A two-qubit certificate]\label{to:prop:certificate}
Two qubit carriers $A$ and $B$ with zero Hamiltonian, which are also the output carriers, arrive in turn, and the target $\mathcal P$ removes the coherences of the inputs and sends $|ab\rangle$ to $(I\otimes X^b)(|00\rangle+(-1)^a|11\rangle)/\sqrt2$. Energies are in units of $\kB T\ln2$ and entropies are in bits.
\begin{enumerate}[label=(\alph*), leftmargin=*]
\item \emph{Held.} If $A$ may be held until $B$ arrives, $\mathcal P$ is implemented exactly with no work and no heat: two controlled-$Z$ gates driven by maximally mixed coins, a Hadamard and a CNOT.
\item \emph{What is measured.} Feed $A=|0\rangle$ and, after $A$ has left, $B$ as half of a maximally entangled pair whose other half $Q$ the observer keeps. The leakage $e$ out of the two-dimensional code of ideal outputs on this probe is fixed by three local Pauli settings, and two bound it: $e\le\Pr[X_AX_B=-1]+\Pr[Z_QZ_AZ_B=-1]$. At strategy error $\varepsilon$, $e\le\varepsilon$.
\item \emph{Work.} Let an early-release implementation have a controller, Gibbs reservoirs at temperature $T$ and a sharp battery that start independent, and on this probe let the changes of the controller's energy and entropy, the battery's entropy, and the battery's mean deviation from its nominal debit total at most $P$. Then its nominal work per episode is at least $1-h_2(e)-e\log_23-P$.
\item \emph{Heat.} No entangled probe is required. For independent probe states, let the device's parts and a monitored bath at temperature $T$ start independent, the bath in its Gibbs state. The net heat into that bath per episode is at least $S(\rho_A)+S(\rho_B)-H_4(p)-\Delta S_C-\Delta S_W-\Delta S_{\rm other}$, where $p$ is the distribution of the two stabilizers above, $H_4(p)\le h_2(e)+e\log_23$, and $\Delta S_C$, $\Delta S_W$ and $\Delta S_{\rm other}$ are the entropy changes of the controller, the battery and every other sink. This form needs no refund of the controller's energy and no sharp battery.
\end{enumerate}
\end{proposition}

The proof is in Appendix~\ref{to:app:certificate}. With $P=0.025$, the work bound of (c) is more than $0.87$ at $\varepsilon=10^{-2}$ and more than $0.05$ up to $\varepsilon\approx0.169$, against $0$ for the held process, and it is tight as $\varepsilon$ and $P$ vanish. For the heat form, let the probe of $B$ be a uniformly random classical bit that the observer records, drawn fresh in every episode, independently of the device's state. A device with memory that knew the setting could tailor its output to it, so the observer chooses the setting of each episode at random, independently of the device's history, and keeps it hidden until both outputs of the episode have been released; every setting then samples the same average conditional emitted state. Take $150{,}000$ episodes, fixed in advance, about $50{,}000$ per setting. With $5\%$ observed leakage, the Hoeffding--Azuma inequality~\cite{Hoeffding63,Azuma67}, which holds when the device carries memory across episodes, keeps the run's average leakage below $7\%$ except with probability below $3\times10^{-4}$. At an average leakage of $7\%$, with a readout calibrated for every input state, a preparation entropy of $0.99$ bits and $0.08$ bits of entropy allowed to the rest of the hardware, the expected net heat into the monitored bath exceeds $0.42\,\kB T\ln2$ per episode for any early-release device, while the held process can be run with no heat at all. Reading the run's average as the expected leakage assumes that the episodes repeat independently, as in the protocol of Section~\ref{sec:protocol} (Appendix~\ref{to:app:certificate}). The measured leakage certifies a lower bound on work or heat on this probe; it is not a reconstruction of the implemented strategy. The bound covers every implementation; only the early-release construction that would attain it needs reservoirs far too large to build. The open step toward an experiment is a fully accounted reset and control interface, so that the controller and battery budgets of a real device can be certified.

\paragraph{A classical witness.} Order costs work for classical processes as well (Appendix~\ref{app:module}), and the smallest case of Proposition~\ref{to:prop:rank} gives a witness without coherence. A first bit $a$ is returned on its carrier; only then does a second bit $b$ arrive, to be returned as $a\oplus b$ on its own carrier. Held, this is a CNOT and costs nothing. Released early, every device, with any quantum memory and any memory across episodes, must dispose of at least $1-h_2(\delta)$ bits of entropy per episode when $a$ and $b$ are fresh uniform bits, independent of each other, of earlier episodes and of the device's state when each arrives, where $\delta$ is the probability that the first output disagrees with the second output decoded with $b$: the entropy of the device as a whole, its reservoirs included, rises by that much, and a device that restores its memory and controls passes it to its reservoirs and other sinks. With errors $e_1$ and $e_2$ on the two outputs, $\delta\le e_1+e_2$, so for $e_1+e_2\le1/2$ at least $1-h_2(e_1+e_2)$ bits must be disposed of, and this is attained exactly (Appendix~\ref{to:app:rank}). Measuring the disagreement needs no record of $a$. Feedback traps for colloidal particles already estimate the work of one-bit erasure from trajectories to a few hundredths of $\kB T$~\cite{Jun14}; what is missing, as for the quantum certificate, is a fully accounted interface for the device's memory and controls.

% =====================================================================
